费用流在残量网络上反复寻找一条最短增广路。
反向边费用为原边费用的相反数，因此允许撤销之前的流。

建边用 \texttt{add(u,v,cap,cost)}；\texttt{flow(s,t)} 返回 \texttt{\{flow,cost\}}。

\textbf{编号约定：}顶点统一使用 $1$-based 编号 $1,2,\ldots,n$；残量边位置仅在内部使用，为 $0$-based 下标。

两个版本均命名为 \texttt{MCMF}，独立复制其中一个使用，要求 $s\ne t$、容量非负，初始正容量边构成的网络无负费用环。
费用、距离、势能及计算中间值需在 \texttt{long long} 范围内，有限距离需小于对应代码的无穷大常量。
第一版容量与总流量使用 \texttt{int}，建边要求 $u\ne v$；第二版容量与总流量使用 \texttt{long long}，总流量需小于 \texttt{inf}。

第一版要求初始正容量边费用非负，势能 $h$ 初始为零，每轮均用 Dijkstra 求最短路。
第二版允许负费用边但无负费用环，仅在开始时用 SPFA 求源点最短路作为势能，之后每轮也用 Dijkstra。
重赋权为 $w'(u,v)=w(u,v)+h_u-h_v$，源点可达部分的残量边重赋权非负。
每轮 Dijkstra 完整结束后，仅对可达点更新 $h_v\gets h_v+d_v$，本轮单位流费用为 $h_t-h_s$。
即使初始费用非负，增广后的反向边也可能为负，因此不能省略势能直接对原费用跑 Dijkstra。
第一版的势能在 \texttt{flow} 内初始化，建图后只调用一次；其中 $h_s=0$，故直接用 $h_t$ 计算本轮单位流费用。

第二版支持 \texttt{flow(s,t,K)}，返回最多 $K$ 单位流的最小费用；需要恰好 $K$ 时检查返回流量是否等于 $K$。
